\newpage

\section{Децентрализованный блочно-координатный Frank--Wolfe}
\label{sec:dbcfw}

В этой главе рассматривается Decentralized Block-Coordinate Frank--Wolfe (DBCFW).
Это распределенное обобщение блочно-координатного Frank--Wolfe.
Метод опирается на идеи BCFW и децентрализованного FW~\cite{lacoste2013block,wai2017decentralized,vedernikov2023decentralized}.
Метод объединяет три идеи:
\begin{enumerate}
    \item децентрализованный консенсус между вычислительными узлами;
    \item отслеживание градиента для приближения среднего градиента глобальной функции;
    \item стохастический выбор блоков и суженный линейный оракул на каждом узле.
\end{enumerate}
Такой алгоритм естественно продолжает BCFW/SOFW: если в транспортной задаче блоками являются OD-пары или источники, то в общей распределенной постановке блоками являются компоненты декартова произведения допустимого множества.

\subsection{Мотивация и место DBCFW среди FW-методов}

Классический Frank--Wolfe хорошо подходит для задач с простым линейным оракулом и дорогой
проекцией.
Однако в современных задачах машинного обучения часто появляются еще две сложности.
Во-первых, данные распределены между несколькими узлами и не могут быть собраны в одном месте.
Во-вторых, само допустимое множество имеет блочную структуру, поэтому полный линейный оракул
избыточен на каждой итерации.

Decentralized Frank--Wolfe решает первую проблему: каждый агент работает со своей локальной
функцией $f_i$, а согласование происходит через коммуникационный граф.
Block-Coordinate Frank--Wolfe решает вторую проблему: на итерации обновляется только часть
блоков допустимого множества.
DBCFW объединяет эти две идеи.
В результате каждый агент не только использует локальную информацию и общается с соседями,
но и решает линейный оракул только на выбранных блоках.

Такая постановка особенно естественна для задач вида
\[
    D=D_1\times\dots\times D_n,
\]
где каждый блок $D_k$ имеет собственный дешевый линейный оракул.
Если полный оракул по $D$ стоит примерно как сумма оракулов по всем блокам, то обновление
батча размера $B\ll n$ уменьшает стоимость итерации примерно в $n/B$ раз.
Цена этого выигрыша состоит в более шумном направлении и необходимости учитывать ошибки
консенсуса между агентами.

\subsection{Связанные направления}

DBCFW опирается на три линии работ.
Первая линия -- децентрализованный FW с градиентным отслеживанием.
В таких методах каждый агент хранит локальную точку $x_t^i$ и локальную оценку глобального
градиента.
Смешивание точек отвечает за консенсус, а смешивание градиентных оценок отвечает за то,
чтобы каждый агент двигался приблизительно по направлению среднего градиента.

Вторая линия -- блочно-координатный FW.
В BCFW допустимое множество является декартовым произведением, и линейный оракул вызывается
только для случайного блока.
Для выпуклых гладких задач это сохраняет порядок $O(1/t)$ по ожидаемой субоптимальности,
но уменьшает стоимость одной итерации.

Третья линия -- стохастический транспортный FW из предыдущей главы.
SOFW показывает, что блочная случайность может быть не только теоретической конструкцией,
но и практическим способом ускорить задачу транспортного назначения.
DBCFW переносит эту идею на более общую распределенную постановку.

\subsection{Постановка распределенной задачи}

Пусть имеется $N$ узлов.
Узел $i$ имеет доступ только к локальной функции $f_i:D\to\R$.
Требуется решить задачу
\begin{equation}
    \min_{x\in D}
    f(x)
    :=
    \frac{1}{N}\sum_{i=1}^{N}f_i(x),
    \label{eq:dbcfw-problem}
\end{equation}
где $D\subset\R^d$ компактно и выпукло.
Предполагается, что множество имеет блочную структуру
\[
    D=D_1\times\dots\times D_n.
\]
Каждый узел поддерживает локальную точку $x_t^i\in D$.
Средняя точка сети равна
\[
    x_{\mathrm{avg},t}
    =
    \frac{1}{N}\sum_{i=1}^{N}x_t^i.
\]
Диаметр допустимого множества обозначается
\[
    \bar\rho=
    \max_{x,y\in D}\norm{x-y}_2.
\]

\subsection{Коммуникационная модель}

На итерации $t$ узлы обмениваются информацией по графу коммуникаций с матрицей смешивания $W_t\in\R^{N\times N}$.
Граф может меняться от итерации к итерации.
Более того, $W_t$ может быть случайной матрицей, если топология коммуникаций или выбор ребер задаются случайным механизмом.
Используются стандартные предположения из децентрализованной оптимизации~\cite{wai2017decentralized,vedernikov2023decentralized}.

\begin{assumption}
\label{ass:dbcfw-mixing}
Для каждой итерации $t$ матрица $W_t$ согласована с текущим графом коммуникаций:
\[
    (W_t)_{ij}=0,
    \qquad
    i\neq j,\quad (i,j)\notin E_t.
\]
Она неотрицательна и дважды стохастична:
\[
    (W_t)_{ij}\geq0,
    \qquad
    W_t\mathbf{1}=\mathbf{1},
    \qquad
    \mathbf{1}^TW_t=\mathbf{1}^T.
\]
Кроме того, существует $\lambda\in(0,1)$ такая, что для всех $z\in\R^N$
\[
    \left\|
    \left(
        W_t-\frac{1}{N}\mathbf{1}\mathbf{1}^T
    \right)z
    \right\|_2
    \leq
    \lambda\norm{z}_2.
\]
Если графы случайны, все эти свойства предполагаются выполненными почти наверное
условно относительно истории до коммуникационного шага.
Иными словами, далее используется одна и та же константа сжатия $\lambda$ для любой
реализации time-varying последовательности матриц $\{W_t\}$.
\end{assumption}

\begin{assumption}
\label{ass:dbcfw-smooth}
Каждая функция $f_i$ выпукла и $L$-гладка на $D$:
\[
    0
    \leq
    f_i(y)-f_i(x)-\inner{\nabla f_i(x)}{y-x}
    \leq
    \frac{L}{2}\norm{y-x}_2^2,
    \qquad
    x,y\in D.
\]
\end{assumption}

Оценки на ошибки консенсуса и градиентного отслеживания в этой главе не постулируются
отдельно.
Они выводятся ниже из предположений о смешивании, гладкости и компактности множества $D$.

\subsection{Алгоритм DBCFW}

В DBCFW каждый узел сначала смешивает локальные точки с соседями:
\[
    \tilde x_t^i=
    \sum_{j=1}^{N}(W_t)_{ij}x_t^j.
\]
Затем он обновляет локальную оценку градиента и снова применяет смешивание:
\[
    g_t^i=
    \tilde g_{t-1}^i
    +
    \nabla f_i(\tilde x_t^i)
    -
    \nabla f_i(\tilde x_{t-1}^i),
\]
\[
    \tilde g_t^i=
    \sum_{j=1}^{N}(W_t)_{ij}g_t^j.
\]
После этого узел выбирает подмножество блоков $I_t^i\subset\{1,\dots,n\}$ размера $B$ и строит суженное множество:
\[
    D_t^i
    =
    \left(\prod_{k\in I_t^i}D_k\right)
    \times
    \left(\prod_{k\notin I_t^i}\{\tilde x_t^i[k]\}\right).
\]
Линейный оракул решается только на выбранных блоках:
\[
    v_t^i
    \in
    \argmin_{v\in D_t^i}
    \inner{\tilde g_t^i}{v}.
\]
Затем выполняется FW-обновление
\[
    x_{t+1}^i
    =
    (1-\gamma_t)\tilde x_t^i+\gamma_t v_t^i.
\]

\begin{algorithm}
\caption{Decentralized Block-Coordinate Frank--Wolfe}
\label{alg:dbcfw}
\begin{algorithmic}
\REQUIRE множество $D=D_1\times\dots\times D_n$, шаги $\gamma_t$, матрицы $W_t$, размер батча $B$
\STATE инициализировать $x_0^i\in D$, $\tilde x_{-1}^i\gets x_0^i$, $\tilde g_{-1}^i\gets\nabla f_i(\tilde x_{-1}^i)$ для всех $i$
\FOR{$t=0,1,2,\ldots$}
\FORALL{$i\in\{1,\ldots,N\}$ параллельно}
\STATE $\tilde x_t^i\gets\sum_{j=1}^{N}(W_t)_{ij}x_t^j$
\STATE $g_t^i\gets\tilde g_{t-1}^i+\nabla f_i(\tilde x_t^i)-\nabla f_i(\tilde x_{t-1}^i)$
\STATE $\tilde g_t^i\gets\sum_{j=1}^{N}(W_t)_{ij}g_t^j$
\STATE выбрать блоки $I_t^i\subset\{1,\ldots,n\}$, $|I_t^i|=B$
\STATE построить $D_t^i=\left(\prod_{k\in I_t^i}D_k\right)\times\left(\prod_{k\notin I_t^i}\{\tilde x_t^i[k]\}\right)$
\STATE $v_t^i\gets\argmin_{v\in D_t^i}\inner{\tilde g_t^i}{v}$
\STATE $x_{t+1}^i\gets(1-\gamma_t)\tilde x_t^i+\gamma_t v_t^i$
\ENDFOR
\ENDFOR
\end{algorithmic}
\end{algorithm}

\subsection{Связь с классическим DFW}

Если каждый узел на каждой итерации выбирает все блоки,
\[
    I_t^i=\{1,\ldots,n\},
\]
то суженное множество совпадает с исходным:
\[
    D_t^i=D.
\]
В этом случае DBCFW превращается в стандартный Decentralized Frank--Wolfe с градиентным отслеживанием:
\[
    v_t^i\in\argmin_{v\in D}\inner{\tilde g_t^i}{v}.
\]
Следовательно, DBCFW является строгим обобщением децентрализованного FW: он сохраняет консенсус и отслеживание градиента, но заменяет полный линейный оракул блочно-суженным оракулом.

\begin{proposition}
\label{prop:dbcfw-to-dfw}
Если для всех узлов $i$ и всех итераций $t$ выбрано $I_t^i=\{1,\ldots,n\}$,
то алгоритм DBCFW совпадает с децентрализованным Frank--Wolfe с градиентным
отслеживанием.
\end{proposition}

\begin{proof}
При полном выборе блоков имеем
\[
    D_t^i
    =
    \left(\prod_{k=1}^{n}D_k\right)=D.
\]
Следовательно, строка линейного оракула в алгоритме DBCFW принимает вид
\[
    v_t^i
    \in
    \argmin_{v\in D}\inner{\tilde g_t^i}{v}.
\]
Все остальные строки алгоритма не зависят от размера батча:
смешивание точек
\[
    \tilde x_t^i=\sum_{j=1}^{N}(W_t)_{ij}x_t^j
\]
и градиентное отслеживание
\[
    g_t^i=
    \tilde g_{t-1}^i+
    \nabla f_i(\tilde x_t^i)-
    \nabla f_i(\tilde x_{t-1}^i)
\]
остаются теми же.
Значит, единственное отличие DBCFW от DFW исчезает, и полученная схема совпадает с
полным децентрализованным FW.
\end{proof}

\subsection{Полный DFW как предельный случай}

Для ясности запишем соответствующую полную схему.
Каждый агент сначала усредняет точки с соседями, затем обновляет оценку градиента,
после чего решает полный линейный оракул.

\begin{algorithm}
\caption{Decentralized Frank--Wolfe с градиентным отслеживанием}
\label{alg:dfw-full}
\begin{algorithmic}
\REQUIRE множество $D$, шаги $\gamma_t$, матрицы смешивания $W_t$
\STATE инициализировать $x_0^i\in D$ и $g_0^i=\nabla f_i(x_0^i)$
\FOR{$t=0,1,2,\ldots$}
\FORALL{$i=1,\ldots,N$ параллельно}
\STATE $\tilde x_t^i\gets\sum_{j=1}^{N}(W_t)_{ij}x_t^j$
\STATE $\tilde g_t^i\gets\sum_{j=1}^{N}(W_t)_{ij}g_t^j$
\STATE $v_t^i\gets\argmin_{v\in D}\inner{\tilde g_t^i}{v}$
\STATE $x_{t+1}^i\gets(1-\gamma_t)\tilde x_t^i+\gamma_tv_t^i$
\STATE $g_{t+1}^i\gets \tilde g_t^i+\nabla f_i(x_{t+1}^i)-\nabla f_i(x_t^i)$
\ENDFOR
\ENDFOR
\end{algorithmic}
\end{algorithm}

DBCFW заменяет только строку полного оракула.
Вместо $D$ используется $D_t^i$, где невыбранные блоки зафиксированы в точке
$\tilde x_t^i$.
Поэтому DBCFW можно рассматривать как DFW с активным подмножеством координат.

\subsection{Связь с SOFW}

SOFW из главы~\ref{sec:scfw} можно рассматривать как активный блочный FW, где агентами являются OD-пары.
Пусть
\[
    x=(x^w)_{w\in OD},
    \qquad
    \cX=\prod_{w\in OD}\cX^w.
\]
Если на итерации выбрано множество активных корреспонденций $\cI_k\subset OD$, то допустимое множество принимает вид
\[
    \cX_k=
    \left(\prod_{w\in\cI_k}\cX^w\right)
    \times
    \left(\prod_{w\notin\cI_k}\{x_k^w\}\right).
\]
Это полностью совпадает с конструкцией $D_t^i$ в DBCFW, если блоки $D_k$ отождествить с блоками $\cX^w$.
Таким образом,
\[
    \text{агент}
    \longleftrightarrow
    \text{OD-пара},
    \qquad
    \text{активные агенты}
    \longleftrightarrow
    \cI_k,
\]
а локальный оракул агента $w$ имеет вид
\[
    \argmin_{u^w\in\cX^w}
    \inner{\tau(f_k)}{\Theta_wu^w}.
\]
В транспортной задаче этот оракул является поиском кратчайшего пути для выбранной корреспонденции.

\subsection{Отличие SOFW от обычного DFW}

Несмотря на похожую блочную форму, SOFW и DFW имеют разную информационную структуру.
В DFW каждый агент хранит свою копию полного вектора $x$ и пытается согласовать ее с копиями
других агентов.
Случайность, если она есть, обычно связана с коммуникациями или стохастическими градиентами.
В SOFW же нет нескольких копий полного решения.
Есть один глобальный поток, разложенный по источникам или OD-парам, а случайность возникает
из выбора активных блоков.

Это различие важно для алгоритмической интерпретации.
SOFW можно записать как распределенный active-set FW:
каждая OD-пара $w$ отвечает за свой блок $\cX^w$, а на итерации активируется только
подмножество таких ``агентов''.
Но эти агенты не решают независимые локальные задачи с собственными функциями $f_i$.
Все они используют один и тот же глобальный градиент
\[
    \tau(f_k)=\nabla\Psi(f_k),
\]
который зависит от суммарного потока по всем ребрам.

DBCFW находится между этими двумя картинами.
Как DFW, он имеет несколько вычислительных узлов с локальными функциями.
Как SOFW, он использует активные блоки допустимого множества.
Поэтому DBCFW можно рассматривать как общий язык, в котором:
\[
    \text{децентрализация}
    \quad\text{и}\quad
    \text{блочная стохастичность}
\]
описаны одновременно.

\subsection{Кривизна и оценка сходимости}

Для анализа вводятся блочные константы кривизны.
Пусть $D=D_1\times\dots\times D_n$.
Для блока $k$ и точки $x\in D$ определим точку $y$:
\[
    y[j]=
    \begin{cases}
        x[k]+\gamma(s[k]-x[k]), & j=k,\\
        x[j], & j\neq k,
    \end{cases}
    \qquad
    \gamma\in[0,1].
\]
Блочная константа кривизны равна
\[
    C_k
    =
    \sup
    \frac{2}{\gamma^2}
    \left[
        f(y)-f(x)-\inner{\nabla f(x)}{y-x}
    \right].
\]
Глобальная константа кривизны
\[
    C
    =
    \sup
    \frac{2}{\gamma^2}
    \left[
        f(x+\gamma(s-x))
        -
        f(x)
        -
        \gamma\inner{\nabla f(x)}{s-x}
    \right]
\]
удовлетворяет
\[
    C\leq\sum_{k=1}^{n}C_k,
    \qquad
    C\leq L\bar\rho^2.
\]

\begin{lemma}[Консенсусная ошибка для time-varying графов]
\label{lem:dbcfw-consensus}
Пусть выполнено предположение~\ref{ass:dbcfw-mixing}.
Пусть $\alpha\in(0,1]$, шаги удовлетворяют
$0\leq\gamma_t\leq \Gamma t^{-\alpha}$, $\gamma_t\leq1$,
$\Gamma\geq 1$, и пусть $t_0$ -- такое целое число, что
\begin{equation}
\label{eq:dbcfw-t0}
    \lambda
    \leq
    \left(\frac{t_0}{t_0+1}\right)^\alpha
    \frac{1}{1+t_0^{-\alpha}}.
\end{equation}
Тогда для любой реализации time-varying последовательности матриц смешивания, удовлетворяющей
предположению~\ref{ass:dbcfw-mixing}, выполнено
\begin{equation}
\label{eq:dbcfw-consensus-bound}
    \Delta_p^t
    :=
    \max_{1\leq i\leq N}\norm{\tilde x_t^i-x_{\mathrm{avg},t}}_2
    \leq
    \frac{C_p}{t^\alpha},
    \qquad
    C_p=t_0^\alpha\Gamma\sqrt N\,\bar\rho .
\end{equation}
Если графы случайны, оценка выполняется почти наверное и, следовательно, также после
взятия математического ожидания.
\end{lemma}

\begin{proof}
Доказательство ведется pathwise: фиксируем произвольную реализацию последовательности
$\{W_t\}$, для которой выполнено предположение~\ref{ass:dbcfw-mixing}.
Это важно для time-varying случая: в доказательстве используется не фиксированная матрица,
а только равномерное сжатие с одним и тем же $\lambda$ на каждой итерации.

Введем матрицы отклонений
\[
    E_t=
    \begin{bmatrix}
        (x_t^1-x_{\mathrm{avg},t})^T\\
        \vdots\\
        (x_t^N-x_{\mathrm{avg},t})^T
    \end{bmatrix},
    \qquad
    \tilde E_t=
    \begin{bmatrix}
        (\tilde x_t^1-x_{\mathrm{avg},t})^T\\
        \vdots\\
        (\tilde x_t^N-x_{\mathrm{avg},t})^T
    \end{bmatrix}.
\]
Двойная стохастичность $W_t$ сохраняет среднее:
\[
    \frac1N\sum_{i=1}^N\tilde x_t^i=x_{\mathrm{avg},t}.
\]
Поэтому, применяя предположение~\ref{ass:dbcfw-mixing} покоординатно и затем суммируя по
координатам, получаем фробениусово сжатие
\begin{equation}
\label{eq:dbcfw-matrix-contraction}
    \norm{\tilde E_t}_F
    =
    \left\|
        \left(W_t-\frac1N\mathbf 1\mathbf 1^T\right)E_t
    \right\|_F
    \leq
    \lambda\norm{E_t}_F .
\end{equation}
Это равенство не требует постоянного графа: на итерации $t$ используется именно текущая
матрица $W_t$.

Для $1\leq t\leq t_0$ имеем $\tilde x_t^i,x_{\mathrm{avg},t}\in D$: смешанные точки
остаются в $D$ по неотрицательности строк $W_t$ и выпуклости $D$, а средняя точка
также лежит в $D$.
Поэтому
\[
    \norm{\tilde E_t}_F\leq \sqrt N\,\bar\rho
    \leq
    \frac{C_p}{t^\alpha}.
\]
Предположим теперь, что для некоторого $t\geq t_0$ выполнено
$\norm{\tilde E_t}_F\leq C_p/t^\alpha$.
Из FW-обновления
\[
    x_{t+1}^i=(1-\gamma_t)\tilde x_t^i+\gamma_t v_t^i,
    \qquad
    x_{\mathrm{avg},t+1}=(1-\gamma_t)x_{\mathrm{avg},t}
    +\gamma_t v_{\mathrm{avg},t},
\]
где $v_{\mathrm{avg},t}=N^{-1}\sum_i v_t^i$, следует
\[
    x_{t+1}^i-x_{\mathrm{avg},t+1}
    =
    (1-\gamma_t)(\tilde x_t^i-x_{\mathrm{avg},t})
    +
    \gamma_t(v_t^i-v_{\mathrm{avg},t}).
\]
Так как $v_t^i\in D$, $v_{\mathrm{avg},t}\in D$ по выпуклости $D$, и диаметр $D$
равен $\bar\rho$,
\[
    \norm{E_{t+1}}_F
    \leq
    (1-\gamma_t)\norm{\tilde E_t}_F
    +
    \gamma_t\sqrt N\,\bar\rho
    \leq
    \frac{C_p+\Gamma\sqrt N\,\bar\rho}{t^\alpha}.
\]
Применяя сжатие уже с матрицей следующей коммуникации $W_{t+1}$, получаем
\[
    \norm{\tilde E_{t+1}}_F
    \leq
    \lambda
    \frac{C_p+\Gamma\sqrt N\,\bar\rho}{t^\alpha}
    =
    \lambda(1+t_0^{-\alpha})\frac{C_p}{t^\alpha}.
\]
Условие~\eqref{eq:dbcfw-t0} замыкает индукцию:
\[
    \lambda(1+t_0^{-\alpha})\frac{1}{t^\alpha}
    \leq
    \frac{1}{(t+1)^\alpha}.
\]
Следовательно,
\[
    \norm{\tilde E_{t+1}}_F\leq \frac{C_p}{(t+1)^\alpha}.
\]
Поскольку максимум по агентам не превосходит фробениусову норму, получаем
оценку~\eqref{eq:dbcfw-consensus-bound}.
\end{proof}

\begin{lemma}[Ошибка градиентного отслеживания]
\label{lem:dbcfw-gradient-tracking}
Пусть выполнены предположения~\ref{ass:dbcfw-mixing}--\ref{ass:dbcfw-smooth}, а шаги
удовлетворяют условиям леммы~\ref{lem:dbcfw-consensus}.
Обозначим
\[
    \tilde g_t=\frac1N\sum_{i=1}^N\tilde g_t^i.
\]
Тогда для любой реализации time-varying последовательности матриц смешивания выполнено
\begin{equation}
\label{eq:dbcfw-gt-bound}
    \Delta_d^t
    :=
    \max_{1\leq i\leq N}\norm{\tilde g_t^i-\tilde g_t}_2
    \leq
    \frac{C_g}{t^\alpha},
\end{equation}
где $C_g$ можно взять любой константой, удовлетворяющей
\[
    C_g
    \geq
    \max
    \left\{
        t_0^\alpha
        \max_{1\leq s\leq t_0}
        \left(
            \sum_{i=1}^N\norm{\tilde g_s^i-\tilde g_s}_2^2
        \right)^{1/2},
        \;
        2\sqrt N\,t_0^\alpha(2C_p+\Gamma\bar\rho)L
    \right\}.
\]
Если графы случайны, оценка снова имеет pathwise-смысл и потому выполняется почти наверное.
\end{lemma}

\begin{proof}
Как и в лемме~\ref{lem:dbcfw-consensus}, фиксируем произвольную реализацию матриц
$\{W_t\}$.
Сначала заметим, что среднее значение градиентного трекера телескопически совпадает со
средним локальным градиентом в смешанных точках:
\begin{equation}
\label{eq:dbcfw-tracker-average}
    \tilde g_t
    =
    \frac1N\sum_{i=1}^N\nabla f_i(\tilde x_t^i).
\end{equation}
Действительно, двойная стохастичность дает
$N^{-1}\sum_i\tilde g_t^i=N^{-1}\sum_i g_t^i$, а обновление
\[
    g_t^i=
    \tilde g_{t-1}^i+
    \nabla f_i(\tilde x_t^i)-\nabla f_i(\tilde x_{t-1}^i)
\]
после суммирования по $i$ дает телескопическую рекурсию.
Начальное условие алгоритма задает эту рекурсию корректно.

Введем энергию рассогласования
\[
    S_t=
    \left(
        \sum_{i=1}^N\norm{\tilde g_t^i-\tilde g_t}_2^2
    \right)^{1/2}.
\]
Достаточно доказать $S_t\leq C_g/t^\alpha$.
Для конечного числа начальных итераций $1\leq t\leq t_0$ это следует из определения
$C_g$; соответствующий максимум конечен, поскольку $D$ компактно, а градиенты $f_i$
непрерывны на $D$.

Пусть $t\geq t_0$ и $S_t\leq C_g/t^\alpha$.
Обозначим
\[
    \delta_{t+1}^i
    =
    \nabla f_i(\tilde x_{t+1}^i)-\nabla f_i(\tilde x_t^i),
    \qquad
    \bar\delta_{t+1}
    =
    \frac1N\sum_{i=1}^N\delta_{t+1}^i .
\]
Тогда
\[
    g_{t+1}^i=\tilde g_t^i+\delta_{t+1}^i,
    \qquad
    \tilde g_{t+1}^i=\sum_{j=1}^N(W_{t+1})_{ij}g_{t+1}^j,
\]
и, по двойной стохастичности $W_{t+1}$,
$\tilde g_{t+1}=\tilde g_t+\bar\delta_{t+1}$.
Здесь принципиально используется именно $W_{t+1}$: при time-varying графах матрица
следующей коммуникации не обязана совпадать с $W_t$.

Сжатие~\eqref{eq:dbcfw-matrix-contraction}, примененное к векторам $g_{t+1}^i$, дает
\[
    S_{t+1}
    \leq
    \lambda
    \left(
        \sum_{i=1}^N
        \norm{g_{t+1}^i-\tilde g_{t+1}}_2^2
    \right)^{1/2}.
\]
Так как
\[
    g_{t+1}^i-\tilde g_{t+1}
    =
    (\tilde g_t^i-\tilde g_t)+(\delta_{t+1}^i-\bar\delta_{t+1}),
\]
получаем
\begin{equation}
\label{eq:dbcfw-gt-split}
    \left(
        \sum_{i=1}^N
        \norm{g_{t+1}^i-\tilde g_{t+1}}_2^2
    \right)^{1/2}
    \leq
    S_t+
    \left(
        \sum_{i=1}^N
        \norm{\delta_{t+1}^i-\bar\delta_{t+1}}_2^2
    \right)^{1/2}.
\end{equation}

Оценим второе слагаемое.
Из $L$-гладкости
\[
    \norm{\delta_{t+1}^i}_2
    \leq
    L\norm{\tilde x_{t+1}^i-\tilde x_t^i}_2.
\]
Для time-varying графа
\[
    \tilde x_{t+1}^i
    =
    \sum_{j=1}^N(W_{t+1})_{ij}x_{t+1}^j,
\]
поэтому, используя стохастичность строк $W_{t+1}$,
\[
\begin{aligned}
    \norm{\tilde x_{t+1}^i-\tilde x_t^i}_2
    &\leq
    \sum_{j=1}^N(W_{t+1})_{ij}
    \norm{x_{t+1}^j-\tilde x_t^j}_2
    +
    \sum_{j=1}^N(W_{t+1})_{ij}
    \norm{\tilde x_t^j-\tilde x_t^i}_2                                   \\
    &\leq
    \gamma_t\bar\rho+2\Delta_p^t
    \leq
    \frac{\Gamma\bar\rho+2C_p}{t^\alpha}.
\end{aligned}
\]
Здесь первая сумма контролируется FW-шагом, а вторая -- леммой
о консенсусе~\ref{lem:dbcfw-consensus}.
Следовательно,
\[
    \norm{\delta_{t+1}^i}_2
    \leq
    \frac{(\Gamma\bar\rho+2C_p)L}{t^\alpha}.
\]
Поскольку $\bar\delta_{t+1}$ является средним этих приращений,
\[
    \norm{\delta_{t+1}^i-\bar\delta_{t+1}}_2
    \leq
    2\frac{(\Gamma\bar\rho+2C_p)L}{t^\alpha},
\]
и
\[
    \left(
        \sum_{i=1}^N
        \norm{\delta_{t+1}^i-\bar\delta_{t+1}}_2^2
    \right)^{1/2}
    \leq
    2\sqrt N\frac{(\Gamma\bar\rho+2C_p)L}{t^\alpha}
    \leq
    \frac{C_g}{t_0^\alpha t^\alpha}.
\]
Подставляя это и индукционное предположение в~\eqref{eq:dbcfw-gt-split}, имеем
\[
    S_{t+1}
    \leq
    \lambda
    \left(1+t_0^{-\alpha}\right)
    \frac{C_g}{t^\alpha}.
\]
Условие~\eqref{eq:dbcfw-t0} снова дает
\[
    S_{t+1}\leq\frac{C_g}{(t+1)^\alpha}.
\]
Индукция доказана, а из $\max_i\norm{\tilde g_t^i-\tilde g_t}_2\leq S_t$ следует
оценка~\eqref{eq:dbcfw-gt-bound}.
\end{proof}

\begin{theorem}
\label{thm:dbcfw-rate}
Пусть выполнены предположения~\ref{ass:dbcfw-mixing}--\ref{ass:dbcfw-smooth}.
Пусть число блоков равно $n$, а размер батча удовлетворяет $1\leq B\leq n$,
\[
    \tau=\frac{2n}{B},
    \qquad
    \gamma_t=\frac{2n}{B(t+\tau)}.
\]
Для $\alpha=1$ и $\Gamma=2n/B$ леммы~\ref{lem:dbcfw-consensus}
и~\ref{lem:dbcfw-gradient-tracking} дают
\[
    \Delta_p^t
    =
    \max_{1\leq i\leq N}\norm{\tilde x_t^i-x_{\mathrm{avg},t}}_2
    \leq\frac{C_p}{t},
    \qquad
    \Delta_d^t
    =
    \max_{1\leq i\leq N}\norm{\tilde g_t^i-\tilde g_t}_2
    \leq\frac{C_g}{t}.
\]
Тогда для среднего итерата при $t\geq1$ выполняется оценка
\[
    \E[f(x_{\mathrm{avg},t})-f(x^\ast)]
    \leq
    \frac{\widetilde A_\tau}{t+\tau},
\]
где
\[
    A_\tau=
    3\bar\rho(C_g+LC_p)+\frac{L}{2}\bar\rho^2,
\]
\[
    H_1=\E[f(x_{\mathrm{avg},1})-f(x^\ast)],
    \qquad
    \widetilde A_\tau
    =
    \max
    \left\{
        (\tau+1)H_1,
        \frac{4A_\tau n^2}{B^2}
    \right\}.
\]
\end{theorem}

\subsection{Доказательство оценки DBCFW}

Доказательство теоремы~\ref{thm:dbcfw-rate} можно разбить на три шага.
Первый шаг связывает изменение глобальной функции в средней точке сети с усредненным
направлением агентов.
Из $L$-гладкости имеем
\[
    f(x_{\mathrm{avg},t+1})
    \leq
    f(x_{\mathrm{avg},t})
    +
    \inner{\nabla f(x_{\mathrm{avg},t})}
    {x_{\mathrm{avg},t+1}-x_{\mathrm{avg},t}}
    +
    \frac{L}{2}
    \norm{x_{\mathrm{avg},t+1}-x_{\mathrm{avg},t}}_2^2.
\]
Так как матрицы $W_t$ дважды стохастичны, средняя точка после смешивания равна средней
точке до смешивания:
\[
    \frac{1}{N}\sum_{i=1}^{N}\tilde x_t^i
    =
    \frac{1}{N}\sum_{i=1}^{N}x_t^i
    =
    x_{\mathrm{avg},t}.
\]
После FW-обновления
\[
    x_{t+1}^i=(1-\gamma_t)\tilde x_t^i+\gamma_t v_t^i
\]
получаем
\[
    x_{\mathrm{avg},t+1}
    =
    x_{\mathrm{avg},t}
    +
    \gamma_t
    (v_{\mathrm{avg},t}-x_{\mathrm{avg},t}),
    \qquad
    v_{\mathrm{avg},t}=\frac{1}{N}\sum_{i=1}^{N}v_t^i.
\]
Следовательно,
\[
    f(x_{\mathrm{avg},t+1})
    \leq
    f(x_{\mathrm{avg},t})
    +
    \gamma_t
    \inner{\nabla f(x_{\mathrm{avg},t})}
    {v_{\mathrm{avg},t}-x_{\mathrm{avg},t}}
    +
    \frac{L\gamma_t^2}{2}\bar\rho^2.
\]

Второй шаг оценивает, насколько $v_{\mathrm{avg},t}$ близок к истинному FW-направлению.
Главная техническая трудность состоит в том, что в DBCFW агент $i$ решает не полный
линейный оракул на $D$, а ограниченный оракул на $D_t^i$.
Поэтому нельзя напрямую сравнить $v_t^i$ с произвольной точкой $a\in D$.
Для такого сравнения вводится точка $a_t^i(a)$, в которой выбранные блоки берутся из $a$,
а остальные блоки замораживаются в смешанной точке $\tilde x_t^i$:
\[
    a_t^i(a)\in D_t^i,
    \qquad
    \left(a_t^i(a)\right)[k]
    =
    \begin{cases}
        a[k], & k\in I_t^i,\\
        \tilde x_t^i[k], & k\notin I_t^i.
    \end{cases}
\]
Из оптимальности $v_t^i$ на суженном множестве следует
\[
    \inner{v_t^i}{\tilde g_t^i}
    \leq
    \inner{a_t^i(a)}{\tilde g_t^i},
    \qquad a\in D.
\]
Теперь сравним направление агента с истинным градиентом в средней точке.
Для любого $a\in D$ имеем
\[
\begin{aligned}
    &\inner{v_t^i-x_{\mathrm{avg},t}}{\nabla f(x_{\mathrm{avg},t})}  \\
    &=
    \inner{v_t^i-x_{\mathrm{avg},t}}{\tilde g_t^i}
    +
    \inner{v_t^i-x_{\mathrm{avg},t}}
    {\nabla f(x_{\mathrm{avg},t})-\tilde g_t^i}  \\
    &\leq
    \inner{a_t^i(a)-x_{\mathrm{avg},t}}{\tilde g_t^i}
    +
    \norm{v_t^i-x_{\mathrm{avg},t}}_2
    \norm{\nabla f(x_{\mathrm{avg},t})-\tilde g_t^i}_2  \\
    &\leq
    \inner{a_t^i(a)-x_{\mathrm{avg},t}}{\nabla f(x_{\mathrm{avg},t})}
    +
    2\bar\rho
    \norm{\nabla f(x_{\mathrm{avg},t})-\tilde g_t^i}_2.
\end{aligned}
\]
В последнем переходе используется компактность $D$:
$\norm{v_t^i-x_{\mathrm{avg},t}}_2\leq\bar\rho$ и
$\norm{a_t^i(a)-x_{\mathrm{avg},t}}_2\leq\bar\rho$.

Осталось оценить рассогласование градиентов.
Обозначим среднее значение переменных трекинга через
\[
    \tilde g_t=\frac{1}{N}\sum_{i=1}^{N}\tilde g_t^i.
\]
Двойная стохастичность $W_t$ и телескопическая форма градиентного трекинга дают тождество
\[
    \tilde g_t
    =
    \frac{1}{N}\sum_{i=1}^{N}\nabla f_i(\tilde x_t^i).
\]
Поэтому
\[
\begin{aligned}
    \norm{\nabla f(x_{\mathrm{avg},t})-\tilde g_t^i}_2
    &\leq
    \norm{\nabla f(x_{\mathrm{avg},t})-\tilde g_t}_2
    +
    \norm{\tilde g_t-\tilde g_t^i}_2                                      \\
    &\leq
    \frac{1}{N}\sum_{j=1}^{N}
    \norm{\nabla f_j(x_{\mathrm{avg},t})-\nabla f_j(\tilde x_t^j)}_2
    +
    \Delta_d^t                                                            \\
    &\leq
    L\Delta_p^t+\Delta_d^t.
\end{aligned}
\]
Следовательно, для любого $a\in D$
\[
    \inner{v_t^i-x_{\mathrm{avg},t}}{\nabla f(x_{\mathrm{avg},t})}
    \leq
    \inner{a_t^i(a)-x_{\mathrm{avg},t}}{\nabla f(x_{\mathrm{avg},t})}
    +
    2\bar\rho(\Delta_d^t+L\Delta_p^t).
\]

Третий шаг использует блочную случайность.
Пусть $\mathcal F_t$ -- история до выбора блоков на итерации $t$.
Если $I_t^i$ выбирается равномерно без возвращения и $|I_t^i|=B$, то для любого
фиксированного $a\in D$ и любого вектора $w$ выполняется
\[
    \E\left[
        \inner{a_t^i(a)-x_{\mathrm{avg},t}}{w}
        \;\middle|\;\mathcal F_t
    \right]
    =
    \frac{B}{n}\inner{a-x_{\mathrm{avg},t}}{w}
    +
    \left(1-\frac{B}{n}\right)
    \inner{\tilde x_t^i-x_{\mathrm{avg},t}}{w}.
\]
Усреднение по агентам зануляет второй член, поскольку
$N^{-1}\sum_i\tilde x_t^i=x_{\mathrm{avg},t}$.
Берем $w=\nabla f(x_{\mathrm{avg},t})$ и выбираем
\[
    a=s_t\in
    \argmin_{s\in D}
    \inner{\nabla f(x_{\mathrm{avg},t})}{s}.
\]
Тогда из выпуклости
\[
    \inner{s_t-x_{\mathrm{avg},t}}{\nabla f(x_{\mathrm{avg},t})}
    \leq
    -\left(f(x_{\mathrm{avg},t})-f(x^\ast)\right).
\]
В итоге получаем условную оценку спуска
\[
    \E\left[
    \inner{\nabla f(x_{\mathrm{avg},t})}{v_{\mathrm{avg},t}-x_{\mathrm{avg},t}}
    \;\middle|\; \mathcal F_t\right]
    \leq
    -\frac{B}{n}
    (f(x_{\mathrm{avg},t})-f(x^\ast))
    +
    2\bar\rho(\Delta_d^t+L\Delta_p^t).
\]
После подстановки в гладкостную оценку получается рекурсия
\[
    h_{t+1}
    \leq
    \left(1-\frac{B}{n}\gamma_t\right)h_t
    +
    2\gamma_t\bar\rho(\Delta_d^t+L\Delta_p^t)
    +
    \frac{L}{2}\gamma_t^2\bar\rho^2,
\]
где
\[
    h_t=\E[f(x_{\mathrm{avg},t})-f(x^\ast)].
\]
Используя оценки лемм~\ref{lem:dbcfw-consensus} и~\ref{lem:dbcfw-gradient-tracking},
\[
    \Delta_p^t\leq\frac{C_p}{t},
    \qquad
    \Delta_d^t\leq\frac{C_g}{t},
\]
получаем
\[
    h_{t+1}
    \leq
    \left(1-\frac{B}{n}\gamma_t\right)h_t
    +
    \gamma_t
    \frac{2\bar\rho(C_g+LC_p)}{t}
    +
    \frac{L}{2}\gamma_t^2\bar\rho^2.
\]
При выборе
\[
    \gamma_t=\frac{2n}{B(t+\tau)},
    \qquad
    \tau=\frac{2n}{B},
\]
имеем
\[
    \frac{B}{n}\gamma_t=\frac{2}{t+\tau}.
\]
Кроме того,
\[
    \frac{\gamma_t/t}{\gamma_t^2}
    =
    \frac{B(t+\tau)}{2nt}
    =
    \frac{B}{2n}+\frac{1}{t}
    \leq
    \frac12+1
    =
    \frac32,
\]
где использованы $B\leq n$ и $t\geq1$.
Следовательно, $\gamma_t/t\leq(3/2)\gamma_t^2$ для $t\geq1$.
Значит, предыдущая рекурсия переходит в более удобный вид
\[
    h_{t+1}
    \leq
    \left(1-\frac{B}{n}\gamma_t\right)h_t
    +
    \left[
        3\bar\rho(C_g+LC_p)+\frac{L}{2}\bar\rho^2
    \right]\gamma_t^2.
\]
Обозначая
\[
    A_\tau=
    3\bar\rho(C_g+LC_p)+\frac{L}{2}\bar\rho^2,
\]
и используя $\gamma_t^2=4n^2/[B^2(t+\tau)^2]$, перепишем рекурсию как
\[
    h_{t+1}
    \leq
    \left(1-\frac{2}{t+\tau}\right)h_t
    +
    \frac{4A_\tau n^2}{B^2(t+\tau)^2}.
\]
Докажем индукцией, что
\[
    h_t\leq\frac{\widetilde A_\tau}{t+\tau},
    \qquad
    \widetilde A_\tau
    =
    \max
    \left\{
        (\tau+1)h_1,
        \frac{4A_\tau n^2}{B^2}
    \right\}.
\]
База при $t=1$ следует прямо из определения $\widetilde A_\tau$:
\[
    h_1
    \leq
    \frac{\widetilde A_\tau}{\tau+1}.
\]
Если оценка верна для некоторого $t\geq1$, то
\[
\begin{aligned}
    h_{t+1}
    &\leq
    \left(1-\frac{2}{t+\tau}\right)
    \frac{\widetilde A_\tau}{t+\tau}
    +
    \frac{4A_\tau n^2}{B^2(t+\tau)^2}                                      \\
    &\leq
    \left(1-\frac{2}{t+\tau}\right)
    \frac{\widetilde A_\tau}{t+\tau}
    +
    \frac{\widetilde A_\tau}{(t+\tau)^2}
    =
    \frac{\widetilde A_\tau(t+\tau-1)}{(t+\tau)^2}.
\end{aligned}
\]
Так как
\[
    \frac{t+\tau-1}{(t+\tau)^2}
    \leq
    \frac{1}{t+\tau+1},
\]
получаем
\[
    h_{t+1}\leq\frac{\widetilde A_\tau}{t+\tau+1}.
\]
Индукция замкнута.
Так сохраняется порядок $O(1/t)$, а влияние децентрализации и батчирования переносится
в константу $\widetilde A_\tau$.

Смысл оценки состоит в том, что DBCFW сохраняет порядок $O(1/t)$ для выпуклой гладкой задачи,
а нужное убывание ошибок консенсуса и отслеживания градиента обеспечивается леммами
выше.
Размер батча $B$ влияет на константы: увеличение числа обновляемых блоков делает шаг ближе к полному DFW, но повышает стоимость линейного оракула.

\subsection{Блочный линейный оракул}

Для $D=D_1\times\dots\times D_n$ полный линейный оракул равен
\[
    \argmin_{v\in D}\inner{g}{v}.
\]
Он распадается по блокам:
\[
    v[k]\in\argmin_{u\in D_k}\inner{g[k]}{u}.
\]
DBCFW использует только выбранный набор $I_t^i$:
\[
    v_t^i[k]\in
    \begin{cases}
        \argmin_{u\in D_k}\inner{\tilde g_t^i[k]}{u}, & k\in I_t^i,\\
        \tilde x_t^i[k], & k\notin I_t^i.
    \end{cases}
\]
Таким образом, невыбранные блоки не меняются в направлении оракула.
После FW-шага они также остаются ближе к текущей смешанной точке:
\[
    x_{t+1}^i[k]=\tilde x_t^i[k],
    \qquad k\notin I_t^i,
\]
если записывать $v_t^i[k]=\tilde x_t^i[k]$.

\subsection{Специальные случаи DBCFW}

DBCFW содержит несколько известных схем как частные случаи.
Если $N=1$, то нет децентрализации и коммуникаций.
Алгоритм превращается в стохастический блочно-координатный FW:
\[
    x_{t+1}=(1-\gamma_t)x_t+\gamma_tv_t,
\]
где $v_t$ отличается от $x_t$ только на выбранных блоках.

Если $B=n$, то каждый агент выбирает все блоки.
Алгоритм превращается в DFW с градиентным отслеживанием.

Если одновременно $N=1$ и $B=n$, то DBCFW становится классическим FW.
Именно поэтому DBCFW удобно рассматривать как общий каркас:
\[
    \text{FW}
    \subset
    \text{BCFW}
    \subset
    \text{DBCFW},
    \qquad
    \text{DFW}
    \subset
    \text{DBCFW}.
\]

\subsection{Интерпретация констант в оценке}

В теореме~\ref{thm:dbcfw-rate} константа
\[
    A_\tau=
    3\bar\rho(C_g+LC_p)+\frac{L}{2}\bar\rho^2
\]
содержит два типа ошибок.
Слагаемое
\[
    3\bar\rho(C_g+LC_p)
\]
отвечает за несовпадение локальных направлений с направлением, которое использовал бы
централизованный алгоритм.
Здесь $C_g$ связан с ошибкой градиентного отслеживания, а $LC_p$ возникает из-за того,
что локальные точки отличаются от средней точки.
Слагаемое
\[
    \frac{L}{2}\bar\rho^2
\]
является обычной гладкостной ценой FW-шага.

Параметр
\[
    \tau=\frac{2n}{B}
\]
увеличивается при малом батче.
Если $B$ мал, шаг должен быть осторожнее, потому что один оракул покрывает меньшую часть
координат.
Если $B=n$, то $\tau=2$, и схема ближе к полному DFW.

\subsection{Итоги раздела}

DBCFW формулируется как общий метод для распределенной проекционно-свободной
оптимизации на декартовых множествах.
Главный алгоритмический объект -- локальный блочный линейный оракул, который вызывается
после смешивания точек и градиентного отслеживания.
Главный теоретический результат -- сохранение порядка $O(1/t)$: леммы о консенсусе и
градиентном отслеживании дают нужное убывание ошибок, а теорема использует эти оценки
в рекурсии спуска.
Главная связь с предыдущими главами состоит в том, что стохастический выбор блоков из SOFW
становится частью распределенного FW-каркаса.


\subsection{Выводы по DBCFW}

DBCFW завершает хронологию развития методов в этой работе.
Сначала предыдущие направления использовались для ускорения FW в одной централизованной задаче.
Затем CFW и NFW были перенесены на машинное обучение.
После этого блочная структура позволила стохастически уменьшить стоимость линейного оракула.
DBCFW объединяет блочную экономию вычислений с распределенной архитектурой, где данные и функции разделены между узлами.

Алгоритм особенно полезен в задачах, где:
\begin{enumerate}
    \item допустимое множество имеет декартову блочную структуру;
    \item линейный оракул по полному множеству слишком дорог по сравнению с блочным оракулом;
    \item данные или функции распределены между несколькими агентами;
    \item коммуникации между агентами ограничены сетевым графом.
\end{enumerate}
Таким образом, DBCFW является естественным кандидатом для крупномасштабных распределенных задач машинного обучения и сетевой оптимизации.
